% /root/monodromy, 2026-11-19. New candidate; independent audit pending.
\subsection{The first-image sheaf and the remaining even-degree gate}
\label{subsec:ss-even-image-sheaf-candidate}
\status{Candidate. The statements below isolate a concrete determinant
divisor on projective space. Its geometric irreducibility and reducedness
are not asserted. The conditional group theorem does not close the
uniform even-degree image problem.}

Continue with $E:y^2+y=x^3$, $X=E^2$ and
$L=\mathcal O_E(3O)\boxtimes\mathcal O_E(3O)$ over an
algebraically closed field of characteristic two. Let $d\ge2$ be
even, $m=3d$, $N=L^d$, $W=H^0(N)$ and $b=m^2=9d^2$.
Write $P=\mathbf P(W)$, $\mathcal B=W^*\otimes\mathcal O_P(-1)$,
and let $\delta_1,\delta_2$ be the canonical lifted invariant
derivations on $N=F_X^*N_{m/2}$.

\begin{lemma}[A negative reflexive extension of the unit image]
\label{lem:ss-even-image-sheaf-candidate}
The space $W$ is a free module of rank $b/4$ over
\[
 \mathcal A=k[\delta_1,\delta_2]/(\delta_1^2,\delta_2^2).
\]
Define reflexive sheaves on $P$ by
\[
 \mathcal I=\ker\bigl(\mathcal B\xrightarrow{(s,\delta_1s,\delta_2s)}
                                  \mathcal O_P^3\bigr),\qquad
 \mathcal K=\ker\bigl(\mathcal B\xrightarrow{
 (s,\delta_1s,\delta_2s,\delta_1\delta_2s)}\mathcal O_P^4\bigr).
\]
Let $Z_0\subset P$ be the linear locus $\delta_1\delta_2s=0$.
It has codimension $b/4$. On $P^\circ=P\setminus Z_0$ these are
vector bundles of ranks $b-3,b-4$ and fit into a Frobenius-equivariant
exact sequence
\[
 0\longrightarrow\mathcal K\longrightarrow\mathcal I
        \xrightarrow{\lambda}\mathcal O_{P^\circ}\longrightarrow0,
 \qquad \lambda=\delta_1\delta_2s,
\]
whose quotient has the constant unit Frobenius. Moreover
\[
 \det\mathcal I=\det\mathcal K=\mathcal O_P(-b),\qquad
 H^0(P,\mathcal I)=H^0(P,\mathcal K)=0.
\]
The full coherent Frobenius maps into $\mathcal I$; its restriction
there is generically bijective and has the required relative unit
Lang system. Its determinant, and the determinant on $\mathcal K$,
give the same nonzero homogeneous polynomial
$h_d^{\mathrm u}\in H^0(P,\mathcal O_P(b))$, up to a common
nonzero constant.
\end{lemma}

\begin{proof}
On either elliptic factor the proof of
Theorem~\ref{thm:ss-product-exact-covectors} gives
\[
 F_*\mathcal O_E(mO)=U\otimes\mathcal O_{E^{(2)}}((m/2)O),
 \qquad 0\to\mathcal O\to U\to\mathcal O\to0
\]
with a nonsplit extension, and $\delta$ is a nonzero multiple of
its quotient-followed-by-inclusion map. Twist this sequence by
the positive-degree line bundle $\mathcal O((m/2)O)$.
Its first cohomology vanishes, so on global sections the kernel
and image of $\delta$ agree and both have dimension $m/2$.
Thus $H^0(\mathcal O_E(mO))$ is a free module of rank $m/2$
over $k[\delta]/\delta^2$. Tensoring the two factors proves
the asserted free $\mathcal A$-module structure.

The operator $\delta_1\delta_2$ has rank $b/4$ on this free
module, proving the codimension of $Z_0$.
If $\delta_1\delta_2s\ne0$, the four vectors
\[
                  s,\quad\delta_1s,\quad\delta_2s,
                       \quad\delta_1\delta_2s
\]
are independent. Apply $\delta_1\delta_2$, then $\delta_1$,
then $\delta_2$ to any relation to kill its first three
coefficients; its last coefficient is then zero as well.
This proves all bundle and surjectivity assertions on $P^\circ$.
More precisely the first three vectors lose rank only when
$(a\delta_1+b'\delta_2)s=0$ for some $[a:b']\in\mathbf P^1$.
Every such nonzero nilpotent operator has rank $b/2$.
Its projective kernel incidence has dimension $b/2$, with
generically unique $[a:b']$, so this smaller rank-drop locus
has codimension $b/2-1$ and lies inside $Z_0$.

Each displayed kernel sheaf is reflexive: it is the kernel of
a vector bundle mapping to a torsion-free image sheaf, so the
depth lemma gives the depth-two condition. Since the evaluation
maps are surjective outside a locus of codimension at least two,
their determinant lines on that open extend as
$\mathcal O_P(-b)$ over all of $P$.
Their inclusions in $\mathcal B$ give the claimed vanishing of
global sections.

Let $M(s)$ denote the root-removed Cartier multiplication matrix,
so full Frobenius has matrix $M(s)^t$ with the usual Frobenius
twists. The previously proved identities are
\[
 M(s)s=M(s)\delta_1s=M(s)\delta_2s=0,
 \qquad M(s)\delta_1\delta_2s=(\delta_1\delta_2s)^{[2]}.
\]
They may also be read off in a pullback frame from
$s=a^2+x_1b_1^2+x_2b_2^2+x_1x_2c^2$ by taking the
$x_1x_2$ coefficient. The first three identities put the full
Frobenius image in $\mathcal I$; the fourth says that $\lambda$
is a morphism to the constant unit-Frobenius line and that
$\mathcal K$ is Frobenius-stable. The exact even-degree stable
rank is $b-3$, so Frobenius on $\mathcal I$ is generically
bijective. Every unit vector of the full coherent bundle lies
in its Frobenius image, hence in $\mathcal I$; consequently the
two unit Lang systems agree on their common dense open.

The determinant on $\mathcal I$ is a regular section of
$\det\mathcal I\otimes(F^*\det\mathcal I)^{-1}
=\mathcal O_P(b)$ away from $Z_0$.
Normality extends it across this codimension-at-least-nine locus.
The quotient in the displayed exact sequence is Frobenius-unit,
so the determinant on $\mathcal K$ agrees with it. Generic
bijectivity proves that this homogeneous polynomial is nonzero.
\end{proof}

\begin{lemma}[An even-degree point with stable loss exactly one]
\label{lem:ss-even-one-loss-point-candidate}
For every even $d\ge2$, the divisor $h_d^{\mathrm u}=0$ has a
point in $P^\circ$ at which the Frobenius-stable rank of
$\mathcal I$ is $b-4$.
\end{lemma}

\begin{proof}
Choose a smooth sharp odd-degree curve $D\in|L^{d-1}|$.
At a general point $P_0\in D$, both coordinate tangent components
are nonzero: each projection $D\to E$ has odd degree $3(d-1)$
and is therefore separable. Choose $P_0$ also away from the two
elliptic flex sets. The codimension-two linear system of
$C\in|L|$ containing $P_0$ and tangent to $D$ there is a
projective six-space.

A generic member of this six-space is globally smooth. At $P_0$
the remaining normal derivative is free. At another point $Q$,
its singularity conditions together with the two imposed
conditions have rank five off the exceptional joint-jet locus,
and at least four on that locus. The latter has dimension at
most one, by the product-cubic first-jet calculation. Thus the
moving singular-member incidence has dimension at most five
inside the indicated six-space. Second jets at $P_0$ are
separated, so the contact can be required to have exactly order
two. The two curve-jet conditions at another point of $D$ have
combined rank four except at finitely many points and at least
three there. The same incidence count excludes every second
double intersection, as in the node-cover proof of
Theorem~\ref{thm:ss-uniform-GL8-block}.

The sharp-source open also meets this six-space. It cannot lie
in the cubic $Q$, which contains no projective linear space of
dimension greater than three. It does not lie in $\det S=0$,
either: normalize the two factor values at $P_0$ to $e_0$ and
their tangent directions to $e_1$. Then the conditions are
$S_{00}=0$ and $aS_{10}+b'S_{01}=0$, with $a,b'\ne0$.
Choose nonzero $S_{10},S_{01}$ satisfying that equation,
set $S_{22}=1$, and all other entries zero. This is an
invertible matrix. All the nonempty opens just described
therefore intersect in the irreducible six-space.

For the resulting $Y=C\cup D$, the sole nonnodal singularity
is a simple tacnode between the two components.
Each normalized component has defect two, and the tacnode
contributes $\delta-r+1=1$. The normalization formula gives
\[
                 g(Y)-f(Y)=2+2+1=5,
 \qquad f(Y)=b-4.
\]
There are also $18(d-1)-2>0$ ordinary intersection nodes.
At such a node the local equation is a product of two transverse
parameters, and $\delta_1\delta_2s$ has nonzero value: it is the
mixed second coefficient in the invariant tangent coordinates.
Thus $Y\notin Z_0$. On $P^\circ$ the full Frobenius image lies
in $\mathcal I$, so its stable rank equals the stable rank on
$\mathcal I$. The full and relative stable ranks agree, since
the ambient supersingular cohomology is nilpotent. This gives
the asserted point and rank.
\end{proof}

\begin{theorem}[A precise conditional even-degree full-image reduction]
\label{thm:ss-even-full-group-if-unit-Hasse-irreducible-candidate}
Let $k$ be any perfect field of characteristic two, and fix even
$d\ge2$ on the indicated surface and embedding. If the divisor
$h_d^{\mathrm u}=0$ is
geometrically irreducible and reduced, then the entire geometric
and arithmetic images are the full pointwise stabilizer of the
unique nonzero invariant covector:
\[
 G_{\mathrm{geom},d}=G_{\mathrm{arith},d}
   =\mathbf F_2^{\,9d^2-4}\rtimes\GL_{9d^2-4}(\mathbf F_2)
       \subset\GL_{9d^2-3}(\mathbf F_2).
\]
\end{theorem}

\begin{proof}
First work over an algebraic closure. The determinant-allocation
proof of Lemma~\ref{lem:irreducible-hasse-quotients-trivial}
applies to $\mathcal I$. We give the only changes required by
its reflexive extension. Saturation of an invariant subspace in
$\mathcal I$ is reflexive by the depth lemma, and its quotient
is torsion-free. In codimension one all terms are vector bundles,
so their effective Frobenius determinant divisors add to the
reduced irreducible $h_d^{\mathrm u}$ divisor. A zero determinant
divisor makes the corresponding reflexive Frobenius map an
isomorphism everywhere, by reflexive extension across codimension
two. Lemma~\ref{lem:periodic-reflexive-projective-trivial} then
makes that sheaf a trivial unit-Frobenius bundle.
Allocating no divisor to a nonzero subobject would contradict
$H^0(P,\mathcal I)=0$. Thus every proper nonzero invariant
subspace has trivial quotient.

The full invariant-covector space has dimension one, by
Theorem~\ref{thm:ss-product-exact-covectors}.
Consequently the common kernel $K$ of that covector is the
only proper nonzero invariant subspace: any such subspace has
constant quotient of dimension at most one, and the quotient
covector is the unique one. In particular $K$ is irreducible
and the whole representation has no invariant line.

The preceding lemma and irreducibility of the determinant
divisor imply generic stable loss exactly one there, by lower
semicontinuity of the rank of a high Frobenius iterate.
The Lang finite-part and purity argument of
Lemma~\ref{lem:char2-defect-one-inertia} therefore makes the
geometric group generated by transvections. Removing $Z_0$
causes no change to that purity argument: it has codimension
at least nine, so the quotient cover is still unramified at
every codimension-one point of the regular projective space.
Restriction to $K$ is thus an irreducible group generated by
transvections (discard any generators restricting to identity).
The actual natural $\GL_8(\mathbf F_2)$ block lies in $K$ and
excludes invariant alternating forms on it, by
Theorem~\ref{thm:ss-uniform-GL8-block}.
The exact classification
\cite[Theorem~3.6, printed p.~22]{KantorSeress2001}
therefore gives the full image $\GL(K)$.

The kernel of the map from the whole group to $\GL(K)$ lies
in the translation module $K$ and is either zero or all of $K$.
If zero, the subgroup is the graph of a one-cocycle on $\GL(K)$.
The proved elementary vanishing
$H^1(\GL_n(\mathbf F_2),\mathbf F_2^n)=0$ for $n\ge4$
(Lemma~\ref{lem:char2-group-inputs}) makes it conjugate to a
linear complement. That complement has a fixed line outside
$K$, contradicting the preceding determinant-allocation result.
Thus every translation occurs, proving the stated geometric group.
The covector is defined over $\mathbf F_2$, so the arithmetic
image is contained in the same stabilizer and contains the
geometric image. Equality follows, with the same dense-open and
ambient-linear-system passage used in
Theorem~\ref{thm:ss-all-odd-full-groups}.
\end{proof}

\begin{remark}[What remains unproved]
The homogeneous polynomial $h_d^{\mathrm u}$ is now specified
intrinsically, with its degree, vector bundle on a complement of
known codimension, and one genuine stable-loss-one point proved
in the preceding lemmas. Its geometric irreducibility and
reducedness for all even $d$ have not been proved here.
The retained exploratory programs
\path{agents/monodromy/code/ss_even_unit_hasse_pencils.py} and
\path{agents/monodromy/code/ss_even_generic_hasse_pencils.py}
compute exact degree-two quotient determinants on six constant-minor
and four moving-minor pencils. Their outputs retain both repeated
and squarefree pencil factorizations. They do not establish a
uniform geometric assertion or replace the missing hypothesis.
\end{remark}
